\chapter{Фотометрия одной серии}
\label{chap:series-photometry}

Это текущая основная рабочая стадия \program{}. В этой главе описана не только последовательность действий, но и математическая модель, которая фактически реализована в коде версии \manualversion. Формулы ниже следует рассматривать как описание текущей реализации: если общепринятый литературный алгоритм отличается от неё, это отмечено явно.

\section{Последовательность этапов}

Типовая цепочка обработки одной серии имеет вид:

\begin{enumerate}
  \item индексирование и чтение FITS metadata;
  \item проверка WCS;
  \item pre-photometry frame QC;
  \item построение или выбор рабочего каталога;
  \item сопоставление catalog sources с кадрами;
  \item source detection и сопоставление detections с каталогом;
  \item forced aperture photometry;
  \item small-to-large aperture correction;
  \item preliminary frame zero point;
  \item internal ensemble calibration;
  \item object QC и field QC;
  \item optional reference/color calibration;
  \item validation random-error model;
  \item materialization финальных science measurements;
  \item публикация candidate и диагностических продуктов.
\end{enumerate}

Общая логика похожа на классическую CCD aperture/differential photometry: измерить поток в апертуре, оценить локальный фон, перевести поток в инструментальную звёздную величину, после чего убрать покадровые общие сдвиги по ансамблю звёзд. Общие принципы aperture photometry и локальной оценки фона обсуждаются, например, у Howell~\cite{howell1989} и Stetson~\cite{stetson1987}; современный практический пример круговой апертуры и локального annulus background приведён также в документации Photutils~\cite{photutils_aperture}. При этом \program{} не вызывает Photutils в основном forced-photometry kernel: текущая реализация использует собственные NumPy-маски.

\section{Обозначения}

В дальнейшем индекс $i$ обозначает кадр, индекс $j$ --- объект. Для одного кадра используются следующие величины:

\begin{description}
  \item[$I_{pq}$] значение пикселя с целочисленными координатами $(p,q)$, в ADU;
  \item[$t_i$] время экспозиции кадра;
  \item[$g$] gain, $e^-/\mathrm{ADU}$;
  \item[$R$] read noise, $e^-$ на пиксель;
  \item[$F_{ij}$] background-subtracted aperture flux объекта $j$ на кадре $i$, в ADU;
  \item[$m^{\mathrm{inst}}_{ij}$] инструментальная звёздная величина;
  \item[$m^{\mathrm{ref}}_j$] каталожная reference magnitude;
  \item[$Z_i$] preliminary aperture zero point кадра;
  \item[$E_i$] внутренний ensemble frame offset;
  \item[$c_j$] цветовой индекс объекта;
  \item[$c_0$] pivot color, используемый в абсолютном color transform.
\end{description}

\section{Детекционная ветвь}
\label{sec:detection-math}

Forced photometry не требует самостоятельной детекции источника: основная позиция берётся из сопоставления рабочего каталога с WCS кадра. Однако параллельно существует detection branch, используемая как независимое evidence для кандидатов, blend diagnostics и QC.

\subsection{Глобальный фон и порог детекции}

Для всех конечных пикселей кадра вычисляется глобальная медиана
\begin{equation}
  B_{\mathrm{glob}} = \operatorname{median}(I_{pq}),
\end{equation}
а затем median absolute deviation
\begin{equation}
  \operatorname{MAD}_{\mathrm{glob}}
  = \operatorname{median}\left|I_{pq}-B_{\mathrm{glob}}\right|.
\end{equation}

При нормальном распределении робастная оценка масштаба берётся как
\begin{equation}
  \sigma_{\mathrm{glob}} = 1.4826\,\operatorname{MAD}_{\mathrm{glob}}.
\end{equation}
Если эта оценка вырождается, используется обычное стандартное отклонение конечных пикселей.

Порог кандидата:
\begin{equation}
  T = B_{\mathrm{glob}} + k\,\sigma_{\mathrm{glob}},
\end{equation}
где $k$ задаётся параметром \texttt{threshold\_sigma}. Пиксель считается начальным кандидатом, если $I_{pq}>T$ и он не ниже каждого из восьми соседей в окрестности $3\times3$.

Для ранжирования применяется именно peak-based, а не aperture-based S/N:
\begin{equation}
  \mathrm{SNR}_{\mathrm{peak}} =
  \frac{I_{\mathrm{peak}}-B_{\mathrm{glob}}}{\sigma_{\mathrm{glob}}}.
\end{equation}
После сортировки по $\mathrm{SNR}_{\mathrm{peak}}$ применяется минимальное допустимое расстояние между принятыми кандидатами. Оптимизированная реализация использует spatial hash, но сохраняет прежнюю greedy-семантику.

\subsection{Уточнение центра}

В небольшой квадратной области вокруг local maximum берутся только положительные после вычитания глобального фона веса
\begin{equation}
  w_{pq}=\max(I_{pq}-B_{\mathrm{glob}},0).
\end{equation}
Центр определяется как центр масс
\begin{equation}
  x_c=\frac{\sum w_{pq}x_{pq}}{\sum w_{pq}},
  \qquad
  y_c=\frac{\sum w_{pq}y_{pq}}{\sum w_{pq}}.
\end{equation}
Если сумма весов неположительна, остаётся исходная координата local maximum. Эта детекционная центровка не заменяет catalog-forced position в основной ветви фотометрии.

\section{Forced aperture photometry одного объекта}
\label{sec:forced-aperture-math}

Основной численный kernel находится в модуле
\texttt{grandgousier/photometry/forced\_aperture.py}. Для каждого catalog-frame link известна позиция $(x,y)$ в пиксельных координатах массива FITS. В текущем алгоритме измерение выполняется в этой позиции без обязательного повторного centroid fit.

По умолчанию рабочая апертура имеет радиус $r_{\rm ap}=3$ pix, а локальный фон оценивается в кольце $8\le r\le12$ pix. Это параметры, а не фундаментальные константы: instrument profile или CLI может задавать другие значения.

\subsection{Геометрия апертуры}

Пиксель относится к aperture set $A$, если его центр удовлетворяет условию
\begin{equation}
  (p-x)^2+(q-y)^2 \le r_{\rm ap}^2.
\end{equation}
Для sky annulus $S$ используется
\begin{equation}
  r_{\rm in}^2 \le (p-x)^2+(q-y)^2 \le r_{\rm out}^2.
\end{equation}

Таким образом, текущая реализация использует бинарное включение пикселя по положению его центра. Это важно отличать, например, от fractional/exact aperture weighting, доступного в Photutils~\cite{photutils_aperture}. Для временных рядов \program{} приоритетом является одинаковая и воспроизводимая геометрия для всех кадров одной серии.

\subsection{Локальный фон}

Локальный фон на пиксель оценивается медианой annulus:
\begin{equation}
  B_{ij}=\operatorname{median}_{(p,q)\in S} I_{pq}.
\end{equation}

Разброс пикселей annulus:
\begin{equation}
  \sigma_{\rm ann}
  =1.4826\,\operatorname{median}_{(p,q)\in S}|I_{pq}-B_{ij}|.
\end{equation}
Если робастная оценка не конечна или равна нулю, используется стандартное отклонение annulus. Если конечных пикселей в annulus меньше пяти, включается fallback на median/std всего локального cutout и выставляется диагностический flag.

Использование локального annulus вместо единственного глобального background является стандартным подходом aperture photometry и особенно важно при пространственно неоднородном фоне~\cite{howell1989,photutils_aperture}.

\subsection{Поток}

Пусть $N_A$ --- число конечных пикселей внутри рабочей апертуры. Суммарный background-subtracted поток равен
\begin{equation}
  F_{ij}
  = \sum_{(p,q)\in A}I_{pq}-N_A B_{ij}.
  \label{eq:aperture-flux}
\end{equation}

Если $F_{ij}\le0$, измерение сохраняется с соответствующим QC flag, но инструментальная magnitude не вычисляется.

\subsection{Инструментальная звёздная величина}

При $F_{ij}>0$ определяется count rate
\begin{equation}
  f_{ij}=\begin{cases}
  F_{ij}/t_i, & t_i>0,\\
  F_{ij}, & \text{если корректная экспозиция отсутствует},
  \end{cases}
\end{equation}
и инструментальная величина
\begin{equation}
  m^{\rm inst}_{ij}
  =-2.5\log_{10}f_{ij}+Z_0.
  \label{eq:instrumental-mag}
\end{equation}
В основном pipeline $Z_0$ является произвольной инструментальной константой, обычно нулём: физический zero point определяется позднее по reference stars.

Логарифмическая форма звёздной величины и использование counts per second являются обычной практикой CCD photometry~\cite{howell1989}.

\section{Модель случайной ошибки aperture measurement}
\label{sec:random-error-model}

Текущая модель \texttt{physical\_aperture\_v1} хранит вклад разных источников шума раздельно. Все дисперсии ниже выражены в $\mathrm{ADU}^2$.

\subsection{Пуассоновский шум источника}

При gain $g$ в $e^-/\mathrm{ADU}$ дисперсия source shot noise после обратного перевода из электронов в ADU равна
\begin{equation}
  \sigma^2_{\rm src}=\frac{\max(F_{ij},0)}{g}.
\end{equation}
Если gain неизвестен, текущий fallback равен $g=1\ e^-/\mathrm{ADU}$ и это явно отмечается в provenance.

\subsection{Шум фона и read noise}

Из detector parameters получается предсказание дисперсии одного background pixel:
\begin{equation}
  \sigma^2_{\rm sky,det}=\frac{\max(B_{ij},0)}{g},
\end{equation}
\begin{equation}
  \sigma^2_{\rm read}=\left(\frac{R}{g}\right)^2,
\end{equation}
\begin{equation}
  \sigma^2_{\rm bg,det}=\sigma^2_{\rm sky,det}+\sigma^2_{\rm read}.
\end{equation}

Одновременно из самого annulus получается эмпирическая дисперсия
\begin{equation}
  \sigma^2_{\rm bg,emp}=\sigma_{\rm ann}^2.
\end{equation}

При включённом empirical background floor программа выбирает
\begin{equation}
  \sigma^2_{\rm bg}=\max\left(
  \sigma^2_{\rm bg,emp},\sigma^2_{\rm bg,det}
  \right).
\end{equation}
Это не позволяет формальной detector model занизить ошибку, если реальные пиксели annulus демонстрируют больший scatter из-за flat-field residuals, неоднородного неба или других эффектов. Необходимость учитывать одновременно source, sky и detector noise является центральной частью классической CCD equation~\cite{howell1989}.

\subsection{Фон внутри апертуры}

Шум background pixels, попавших внутрь aperture sum:
\begin{equation}
  \sigma^2_{\rm ap,bg}=N_A\sigma^2_{\rm bg}.
\end{equation}

Кроме того, сам $B_{ij}$ является оценённой величиной. Для медианы и приблизительно Gaussian background текущий pipeline использует variance factor $\pi/2$:
\begin{equation}
  \sigma^2_{\widehat B}
  =\frac{\pi}{2}\,\frac{N_A^2}{N_S}\,\sigma^2_{\rm bg},
\end{equation}
где $N_S$ --- число конечных annulus pixels. Множитель $N_A^2$ возникает потому, что ошибка оценки background level вычитается из каждого из $N_A$ aperture pixels совместно, а не независимо.

\subsection{Дополнительный fractional floor}

Если instrument policy задаёт относительный floor $f_{\rm floor}$, он записывается как flux-proportional variance
\begin{equation}
  \sigma^2_{\rm floor}=(\max(F_{ij},0)f_{\rm floor})^2.
\end{equation}
По умолчанию этот член равен нулю.

\subsection{Сцинтилляция}

Если известны диаметр телескопа, высота площадки, airmass и время экспозиции, используется Young-type approximation
\begin{equation}
  s_{\rm scint}
  =0.09\,D_{\rm cm}^{-2/3}X^{1.75}
  \exp\left(-\frac{h}{8000\ \mathrm{m}}\right)
  (2t_i)^{-1/2},
  \label{eq:scintillation}
\end{equation}
где $s_{\rm scint}$ --- относительный RMS. Тогда
\begin{equation}
  \sigma^2_{\rm scint}=(\max(F_{ij},0)s_{\rm scint})^2.
\end{equation}

Это приближённая модель: реальная scintillation зависит от вертикального профиля турбулентности и ветра. Теория и масштабирование scintillation с telescope aperture подробно обсуждаются Dravins et al.~\cite{dravins1998}; поэтому этот член следует трактовать как физически мотивированную оценку, а не как точное описание конкретной ночи.

\subsection{Итоговая ошибка}

Полная случайная дисперсия flux:
\begin{equation}
  \sigma_F^2=
  \sigma^2_{\rm src}
  +\sigma^2_{\rm ap,bg}
  +\sigma^2_{\widehat B}
  +\sigma^2_{\rm floor}
  +\sigma^2_{\rm scint}.
  \label{eq:flux-error-total}
\end{equation}

Для положительного потока magnitude uncertainty получается линейным распространением ошибки через $m=-2.5\log_{10}F$:
\begin{equation}
  \sigma_m=\frac{2.5}{\ln 10}\frac{\sigma_F}{F}
  \simeq1.085736\frac{\sigma_F}{F}.
  \label{eq:mag-error}
\end{equation}

Параллельно для provenance вычисляется legacy-модель
\begin{equation}
  \sigma^2_{F,\rm legacy}
  =\max(F,0)+N_A\sigma_{\rm ann}^2,
\end{equation}
но в нормальном режиме она используется только для A/B diagnostics, а не как основная физическая ошибка.

\section{Что означает ``фотометрия одного кадра''}

Для одного кадра предыдущие вычисления повторяются для всех catalog-frame links. Поэтому результат кадра можно представить как набор
\begin{equation}
  \mathcal{M}_i=
  \left\{
  x_{ij},y_{ij},F_{ij},\sigma_{F,ij},
  m^{\rm inst}_{ij},\sigma_{m,ij},B_{ij},\mathrm{flags}_{ij}
  \right\}_{j=1}^{N_i}.
\end{equation}

Важно, что forced measurement и detection measurement --- разные сущности. Detection branch отвечает на вопрос ``есть ли здесь самостоятельный local maximum и насколько он убедителен'', а forced branch --- ``какой поток измеряется в известной позиции catalog object''. Поэтому слабый объект может иметь корректное forced measurement даже без отдельной detection.

\section{Small-to-large aperture correction}
\label{sec:aperture-correction-math}

Малая рабочая апертура выгодна по S/N, особенно для слабых объектов, но её доля enclosed PSF flux может изменяться по кадру и между кадрами. Идея corrections по growth curve/large aperture является стандартной для aperture photometry~\cite{howell1989}. В \program{} correction строится отдельно для каждого кадра.

\subsection{Измерение calibrators}

Из ярких, несатурированных и достаточно точных measurements выбирается набор calibrators. Каждый calibrator повторно измеряется большой reference aperture. Для него вычисляется
\begin{equation}
  d_{ij}=m^{\rm large}_{ij}-m^{\rm small}_{ij}.
  \label{eq:aperture-delta}
\end{equation}

Перед fit выполняется iterative clipping по $d$: если $\widetilde d$ --- медиана, а $\operatorname{MAD}(d)$ --- median absolute deviation, сохраняются точки
\begin{equation}
  |d_{ij}-\widetilde d|
  \le k_{\rm clip}\,1.4826\,\operatorname{MAD}(d).
\end{equation}

\subsection{Constant model}

Самая простая модель:
\begin{equation}
  C_i(x,y)=c_i,
\end{equation}
где
\begin{equation}
  c_i=\operatorname{median}(d_{ij}).
\end{equation}

\subsection{Spatial plane}

Для учёта положения на detector координаты нормируются в диапазон примерно $[-1,1]$:
\begin{equation}
  x_n=2\frac{x}{W-1}-1,
  \qquad
  y_n=2\frac{y}{H-1}-1.
\end{equation}

Альтернативная модель:
\begin{equation}
  C_i(x,y)=c_i+a_i x_n+b_i y_n.
  \label{eq:aperture-plane}
\end{equation}
Коэффициенты $(c_i,a_i,b_i)$ определяются обычным linear least squares по clipped calibrators.

\subsection{Выбор constant или plane}

Для residuals $r_k$ вычисляются
\begin{equation}
  \mathrm{RMS}=\sqrt{\frac{1}{n}\sum_{k=1}^{n}r_k^2},
\end{equation}
\begin{equation}
  \mathrm{MAD}=\operatorname{median}|r_k-\operatorname{median}(r)|,
\end{equation}
и Bayesian Information Criterion
\begin{equation}
  \mathrm{BIC}=n\ln\left(\frac{\mathrm{SSE}}{n}\right)+k\ln n,
  \qquad
  \mathrm{SSE}=\sum r_k^2,
  \label{eq:bic}
\end{equation}
где $k=1$ для constant model и $k=3$ для plane. Это стандартная форма Schwarz criterion~\cite{schwarz1978}.

Plane принимается только если одновременно выполнены статистические и safety gates. В текущей реализации по умолчанию, среди прочего, требуется
\begin{equation}
  \mathrm{BIC}_{\rm const}-\mathrm{BIC}_{\rm plane}\ge6,
\end{equation}
минимальный spatial peak-to-peak effect, отсутствие ухудшения residual MAD и ограничения на максимальную пространственную/полную correction. Если условия не выполнены, используется constant model или identity fallback.

\subsection{Применение correction}

К каждому small-aperture measurement прибавляется
\begin{equation}
  m^{\rm ap}_{ij}
  =m^{\rm inst}_{ij}+C_i(x_{ij},y_{ij}).
  \label{eq:aperture-apply}
\end{equation}
Эквивалентный пересчёт flux:
\begin{equation}
  F^{\rm ap}_{ij}
  =F_{ij}\,10^{-0.4C_i(x_{ij},y_{ij})}.
\end{equation}

Постоянная часть $c_i$ позднее поглощается frame zero point. Поэтому в per-object random error переносится только неопределённость пространственной части plane. Если $\Sigma_i$ --- covariance matrix коэффициентов $(c_i,a_i,b_i)$, а $(\bar x_n,\bar y_n)$ --- средняя позиция aperture calibrators, используется вектор
\begin{equation}
  \mathbf v=
  \begin{bmatrix}
  0 & x_n-\bar x_n & y_n-\bar y_n
  \end{bmatrix}^{\!T},
\end{equation}
и
\begin{equation}
  \sigma^2_{\rm ap,diff}=\mathbf v^T\Sigma_i\mathbf v.
\end{equation}
Итоговая random magnitude error после aperture correction:
\begin{equation}
  \sigma^2_{m,\rm ap}
  =\sigma_m^2+\sigma^2_{\rm ap,diff}.
\end{equation}
Для constant model этот дополнительный дифференциальный член равен нулю.

\section{Preliminary aperture zero point}
\label{sec:aperture-zp-math}

После aperture correction каждый кадр предварительно привязывается к reference magnitudes. Для пригодной reference star $j$ на кадре $i$ вычисляется
\begin{equation}
  z_{ij}=m^{\rm ref}_j-m^{\rm ap}_{ij}.
  \label{eq:zp-candidate}
\end{equation}

Из выборки исключаются, в частности, saturated measurements и measurements с недопустимой magnitude error. Затем выполняется iterative robust clipping относительно медианы:
\begin{equation}
  |z_{ij}-\widetilde z_i|
  \le k_{\rm clip}\,1.4826\,\operatorname{MAD}(z_i).
\end{equation}

Если после clipping осталось достаточно reference stars и frame-MAD проходит quality gate, zero point равен
\begin{equation}
  Z_i=\operatorname{median}_{j}(z_{ij}).
  \label{eq:frame-zp}
\end{equation}
После чего
\begin{equation}
  m^{\rm zp}_{ij}=m^{\rm ap}_{ij}+Z_i.
  \label{eq:zp-apply}
\end{equation}

Эта стадия называется preliminary, поскольку reference catalog задаёт приблизительную абсолютную шкалу, но ещё не устранены остаточные frame-common offsets и color mismatch.

\section{Ансамблевая дифференциальная фотометрия}
\label{sec:ensemble-math}

Классическая идея ensemble photometry состоит в том, чтобы использовать одновременно множество comparison stars вместо одной пары target/comparison. Это уменьшает чувствительность к случайной проблеме отдельной comparison star и позволяет работать с неоднородными наборами наблюдений. Один из основных классических источников --- Honeycutt~\cite{honeycutt1992}.

Следует, однако, различать общую идею и конкретный алгоритм. Honeycutt описывает глобальное least-squares ensemble solution. Текущая реализация \program{} --- \emph{не} прямое воспроизведение этого least-squares решения; она использует устойчивую итеративную схему на медианах frame residuals.

\subsection{Начальная статистика объектов}

Для каждого объекта вычисляются по всем доступным кадрам:
\begin{equation}
  M_j^{(0)}=\operatorname{median}_{i}m^{\rm zp}_{ij},
\end{equation}
sample RMS
\begin{equation}
  \mathrm{RMS}_j
  =\sqrt{\frac{1}{N_j-1}
  \sum_i\left(m^{\rm zp}_{ij}-\bar m_j\right)^2},
\end{equation}
и
\begin{equation}
  \mathrm{MAD}_j
  =\operatorname{median}_{i}
  \left|m^{\rm zp}_{ij}-M_j^{(0)}\right|.
\end{equation}

Reference ensemble $\mathcal R$ формируется по числу measurements, диапазону magnitude, максимальному RMS/MAD и явным exclusions. Target/candidate objects могут быть исключены из reference ensemble, чтобы собственная переменность кандидата не участвовала в нормировке.

\subsection{Frame offset}

На итерации $q$ для каждого кадра берутся residuals reference stars относительно их текущих object medians:
\begin{equation}
  r_{ij}^{(q)}=m^{\rm zp}_{ij}-M_j^{(q)},
  \qquad j\in\mathcal R.
\end{equation}

Frame-common offset:
\begin{equation}
  E_i^{(q)}=
  \operatorname{median}_{j\in\mathcal R}r_{ij}^{(q)}.
  \label{eq:ensemble-frame-offset}
\end{equation}
Кадр получает offset только при достаточном числе reference stars. Если пользовательский минимум
обозначить $N_{\min}$, то необходимое условие для кадра имеет вид
\begin{equation}
  N_{\mathrm{ref},i} \ge N_{\min}.
  \label{eq:ensemble-min-ref}
\end{equation}
Начиная с 0.5.45.1 обычный field pipeline явно завершает такой случай статусом
\path{FAILED_PRECONDITION_ENSEMBLE}, если conventional ensemble не смог сформировать ни одного
полного frame solution. Это предотвращает вторичное падение downstream QC на несуществующей
calibration generation. Параметр \texttt{process-archive-path --ensemble-min-ref-stars-per-frame}
делает $N_{\min}$ доступным из one-command workflow; его default остаётся 30 для сохранения прежней
семантики.

\subsection{Обновление ensemble}

После вычисления offsets reference-star medians пересчитываются:
\begin{equation}
  M_j^{(q+1)}=
  \operatorname{median}_{i}
  \left(m^{\rm zp}_{ij}-E_i^{(q)}\right),
  \qquad j\in\mathcal R.
  \label{eq:ensemble-object-update}
\end{equation}

Шаги~\eqref{eq:ensemble-frame-offset}--\eqref{eq:ensemble-object-update} повторяются заданное число внутренних итераций. Затем offset последней итерации применяется ко всем объектам кадра:
\begin{equation}
  m^{\rm ens}_{ij}=m^{\rm zp}_{ij}-E_i.
  \label{eq:ensemble-apply}
\end{equation}

Таким образом, $E_i$ описывает общий additive shift кадра. Его scatter по reference stars сохраняется как QC diagnostic. В текущей error semantics неопределённость frame-common ensemble offset не складывается непосредственно с независимой per-epoch random error каждого объекта: frame-common и independent contributions намеренно не смешиваются.

\section{Candidate-excluded ensemble convergence}

Если включён режим исключения кандидатов из ensemble, после первичной калибровки программа строит список candidate/reference exclusions и может повторять ensemble calibration до сходимости. Внешняя итерация отличается от внутренних итераций формул~\eqref{eq:ensemble-frame-offset}--\eqref{eq:ensemble-object-update}: внешняя итерация может менять само множество $\mathcal R$, тогда как внутренняя фиксирует reference membership и уточняет $M_j$ и $E_i$.

Это различие важно при интерпретации provenance: число внутренних ensemble iterations и число candidate-exclusion convergence iterations являются разными параметрами.

\section{Object QC и field QC}

После ensemble calibration программа вычисляет статистики временного ряда каждого объекта. Основные используемые величины имеют стандартный вид:
\begin{equation}
  \widetilde m_j=\operatorname{median}_i m^{\rm ens}_{ij},
\end{equation}
\begin{equation}
  \mathrm{MAD}_j=\operatorname{median}_i
  |m^{\rm ens}_{ij}-\widetilde m_j|,
\end{equation}
\begin{equation}
  \sigma_{\rm robust,j}=1.4826\,\mathrm{MAD}_j,
\end{equation}
а также sample RMS. Эти величины сопоставляются с magnitude, количеством measurements, detection support и crowding/blend diagnostics. Поэтому QC flag --- результат совокупности критериев, а не просто порога по одному RMS.

Field QC агрегирует object/frame statistics и используется как контроль качества серии целиком. Конкретные пороги относятся к policy/configuration и не являются универсальными астрономическими константами.

\section{Абсолютная reference/color calibration}
\label{sec:color-calibration-math}

Ensemble calibration хорошо сохраняет относительную форму light curve, но её шкала ещё может отличаться от выбранной стандартной/reference системы из-за общего offset и несовпадения spectral response. Стандартная photometric calibration поэтому часто содержит zero point и color term; примеры таких преобразований приведены в руководстве AAVSO по CCD photometry~\cite{aavso_ccd}.

В \program{} absolute transform строится на уровне объектов, а не отдельных эпох. Это сделано специально для сохранения формы временного ряда.

\subsection{Object-level residual}

Для каждого calibrator object сначала вычисляется median ensemble magnitude
\begin{equation}
  S_j=\operatorname{median}_i m^{\rm ens}_{ij}.
\end{equation}
Пусть $m^{\rm ref}_j$ --- reference magnitude в выбранном band, а $c_j$ --- reference color. Pivot color определяется робастно по calibrators:
\begin{equation}
  c_0=\operatorname{median}_j c_j.
\end{equation}

Фитируется уравнение
\begin{equation}
  m^{\rm ref}_j-S_j=A+C(c_j-c_0),
  \label{eq:color-transform-fit}
\end{equation}
где $A$ --- общий transform offset, $C$ --- color term.

Коэффициенты находятся ordinary linear least squares с intercept. После каждого fit residuals робастно clip-ятся и fit повторяется. Параллельно строится offset-only model ($C=0$).

\subsection{Выбор модели}

Для offset-only и offset+color models вычисляются residual RMS, MAD и BIC по формуле~\eqref{eq:bic}. Кроме того, для color slope вычисляется significance
\begin{equation}
  S_C=\frac{|C|}{\sigma_C}.
\end{equation}

Color model принимается консервативно: должны одновременно пройти минимальная significance, улучшение BIC и residual-MAD criterion. Если evidence недостаточно, применяется offset-only transform. Такой выбор защищает от добавления слабого color term только потому, что более сложная модель формально имеет ещё один свободный параметр.

\subsection{Применение к light curve}

Для объекта $j$ одна и та же correction применяется ко всем его эпохам:
\begin{equation}
  \Delta^{\rm color}_j=A+C(c_j-c_0),
\end{equation}
\begin{equation}
  m^{\rm col}_{ij}=m^{\rm ens}_{ij}+\Delta^{\rm color}_j.
  \label{eq:color-apply}
\end{equation}

Поскольку $\Delta^{\rm color}_j$ не зависит от $i$, внутриобъектные RMS, MAD и peak-to-peak не должны изменяться. Это проверяется отдельным shape-preservation quality gate.

Если reference color отсутствует, используется $c_j=c_0$, то есть color contribution к nominal correction равен нулю, но uncertainty расширяется.

\subsection{Систематическая ошибка абсолютного преобразования}

Пусть $x_j=c_j-c_0$. Текущая модель systematic uncertainty абсолютного transform имеет вид
\begin{equation}
\begin{split}
  \sigma^2_{\rm sys,j}={}&
  \sigma^2_{\rm model}
  +\sigma_A^2
  +x_j^2\sigma_C^2\\
  &+2x_j\operatorname{Cov}(A,C)
  +C^2\sigma_{c,j}^2.
\end{split}
\label{eq:color-systematic}
\end{equation}

Эта ошибка хранится как \texttt{mag\_err\_systematic} и является в основном коррелированной между эпохами одного объекта. Она принципиально отделяется от \texttt{mag\_err\_random}.

\section{Полная формула science magnitude}
\label{sec:science-composite-formula}

Если все стадии доступны и приняты quality gates, цепочку для одной точки удобно записать одной формулой:
\begin{equation}
\boxed{
  m^{\rm science}_{ij}
  =m^{\rm inst}_{ij}
  +C^{\rm ap}_i(x_{ij},y_{ij})
  +Z_i
  -E_i
  +A
  +C(c_j-c_0)
  }
  \label{eq:science-full}
\end{equation}

Здесь:
\begin{itemize}
  \item $C^{\rm ap}_i$ --- small-to-large aperture correction;
  \item $Z_i$ --- preliminary reference zero point;
  \item $E_i$ --- robust ensemble frame offset;
  \item $A+C(c_j-c_0)$ --- object-constant absolute reference transform.
\end{itemize}

Формула~\eqref{eq:science-full} полезна для понимания физического смысла pipeline, но в БД эти слагаемые намеренно хранятся как отдельные immutable/normalized calibration terms. Это позволяет восстановить science product, проверить provenance и не переписывать десятки тысяч промежуточных measurements при каждой новой calibration generation.

Поле \texttt{local\_zp\_correction} после ensemble stage соответствует покадровой части $Z_i-E_i$. Color transform туда не входит, поскольку он object-level, а не frame-level.

Следует также помнить, что финальный \texttt{flux} остаётся instrumental aperture flux, скорректированным только на aperture correction. Zero point, ensemble offset и color transform применяются к magnitude representation; \program{} не объявляет \texttt{flux} физически откалиброванным flux density.

\section{Семантика ошибок в science layer}

В текущем science product различаются как минимум три уровня:

\begin{description}
  \item[\texttt{mag\_err\_random}] независимая per-epoch uncertainty из aperture error budget плюс дифференциальная неопределённость spatial aperture correction;
  \item[\texttt{mag\_err\_systematic}] object-level uncertainty абсолютного reference/color transform по формуле~\eqref{eq:color-systematic};
  \item[QC/noise-floor diagnostics] дополнительные empirical residual floors и frame-common components, которые могут храниться в provenance/flags и не обязаны автоматически складываться с \texttt{mag\_err\_random}.
\end{description}

Поле \texttt{mag\_err} в настоящее время сохраняется как backward-compatible представление random per-epoch uncertainty. Поэтому при научной работе, требующей covariance-aware анализа, нельзя механически считать все ошибки независимыми и складывать \texttt{mag\_err\_random} с \texttt{mag\_err\_systematic} для каждой эпохи как будто это независимые Gaussian draws.

\section{Что является литературным методом, а что спецификой Grandgousier}

Для корректной интерпретации полезно разделять общие идеи и конкретные инженерные решения:

\begin{itemize}
  \item circular aperture photometry, local sky estimation, CCD noise budget и growth/aperture corrections имеют классические аналоги в литературе~\cite{howell1989,stetson1987};
  \item ensemble differential photometry основана на общей идее множества comparison stars~\cite{honeycutt1992}, но текущая реализация \program{} использует iterative robust frame medians, а не глобальную least-squares систему Honeycutt;
  \item BIC применяется в стандартном смысле Schwarz~\cite{schwarz1978}, но конкретные пороги принятия plane/color models являются policy проекта;
  \item scintillation term использует приближённое масштабирование Young-type; реальные атмосферные условия сложнее~\cite{dravins1998};
  \item color equation является обычным линейным photometric transform, но \program{} специально делает correction object-constant, чтобы не менять форму ensemble light curve.
\end{itemize}

Поэтому ссылки в приложении~\ref{app:references} объясняют научный контекст методов, а точной спецификацией вычислений для воспроизводимости остаются формулы этой главы, provenance в БД и соответствующая версия исходного кода.

\section{BASE и final SCIENCE storage в 0.7}

В clean schema 0.7 таблица \texttt{measurements} содержит только immutable BASE
photometry. Низкоуровневые writers проверяют, что target calibration имеет
\texttt{product\_kind=BASE\_PHOTOMETRY}, \texttt{product\_role=BASE} и
\texttt{materialization\_kind=FULL\_MEASUREMENTS}; derived full-copy rows в
canonical \texttt{measurements} не допускаются.

Промежуточные calibration stages сохраняются normalized terms. Финальный SCIENCE
product хранится в durable narrow \texttt{science\_photometry}, связанной с BASE
через \texttt{base\_measurement\_id}. В ней сохраняются только calibration-changing
скаляры; extraction geometry и identity не дублируются.

Legacy wide SCIENCE materialization, rebuildable cache и historical flag-override
workflow 0.6 не входят в новую schema и не доступны как штатные CLI-команды 0.7.
Frozen production DB 0.6.4.5.1 используется только как reference corpus для
последующего old/new scientific acceptance.

\subsection{Диагностическая политика выбора spatial plane}
\label{sec:aperture-plane-policy}

По умолчанию Grandgousier использует политику \texttt{evidence}: пространственная плоскость
для aperture correction выбирается только при достаточном числе калибраторов, достаточном
улучшении BIC, ненулевой пространственной амплитуде, residual MAD не хуже constant-модели
и прохождении safety-ограничений по максимальной пространственной и полной поправке.

Для контролируемой диагностики инструментальной систематики доступен режим
\texttt{--aperture-correction-plane-policy safety\_only}. В нём статистические предпочтения
между constant и plane не блокируют плоскость: если plane построен и проходит safety-gates,
он используется. При этом общий порог \texttt{--aperture-correction-max-residual-mad}
остаётся активным и может отвергнуть модель целиком. Режим предназначен для A/B-тестов,
а не для безусловного production-применения.

Пример:
\begin{verbatim}
python -m grandgousier.cli process-archive-path `
  --path C:\\master\\gpx\\stage2\\20180907_1704_R `
  --execute `
  --aperture-correction-plane-policy safety_only
\end{verbatim}

При сравнении с baseline следует проверять не только число кандидатов, но и долю кадров
с моделями plane/constant/identity, spatial coefficients, detector-coordinate correlations
кандидатных light curves и common-mode корреляции между кандидатами.

\section{Диагностика PSF по полю кадра}
\label{sec:psf-field-diagnostics}

Начиная с версии 0.5.44.0 \program{} имеет отдельный диагностический слой для анализа формы
изображений звёзд по полю. Он предназначен прежде всего для широкопольных систем и архивных
данных, где качество изображения может зависеть от положения на детекторе вследствие coma,
astigmatism, field curvature, tilt/decenter, defocus, tracking errors или других оптических и
механических эффектов. На этом этапе диагностика \emph{не изменяет} science flux и не переводит
кадр автоматически в BAD-состояние.

Для каждой подходящей source detection используется компактная окрестность источника. После
вычитания локального фона отрицательные веса обнуляются, затем вычисляются centroid и матрица
вторых моментов
\begin{equation}
  \mathbf{M}=
  \frac{1}{\sum_p w_p}
  \sum_p w_p
  \begin{pmatrix}
    (x_p-\bar{x})^2 & (x_p-\bar{x})(y_p-\bar{y})\\
    (x_p-\bar{x})(y_p-\bar{y}) & (y_p-\bar{y})^2
  \end{pmatrix}.
\end{equation}

Если собственные значения матрицы равны $\lambda_{\rm min}$ и $\lambda_{\rm max}$, диагностические
размеры PSF определяются как
\begin{equation}
  {\rm FWHM}_{\rm minor}=2\sqrt{2\ln2}\sqrt{\lambda_{\rm min}},\qquad
  {\rm FWHM}_{\rm major}=2\sqrt{2\ln2}\sqrt{\lambda_{\rm max}},
\end{equation}
а эллиптичность --- как
\begin{equation}
  e=1-\sqrt{\frac{\lambda_{\rm min}}{\lambda_{\rm max}}}.
\end{equation}
Ориентация главной оси хранится отдельно; из-за $180^\circ$-периодичности дополнительно
сохраняются $\cos 2\theta$ и $\sin 2\theta$.

Чтобы параметры разных детекторов можно было сравнивать, координаты нормируются:
\begin{equation}
  x'=2\frac{x}{N_x-1}-1,\qquad
  y'=2\frac{y}{N_y-1}-1.
\end{equation}

Для FWHM, ellipticity и curve-of-growth diagnostic по каждому кадру независимо строятся три
robust spatial models:
\begin{align}
  q(x',y') &= c_0,\\
  q(x',y') &= c_0+c_xx'+c_yy',\\
  q(x',y') &= c_0+c_xx'+c_yy'+c_{xx}x'^2+c_{xy}x'y'+c_{yy}y'^2.
\end{align}
Fit использует iterative MAD clipping. Для каждой модели сохраняются число использованных
звёзд, коэффициенты, residual MAD, BIC и predicted detector-wide peak-to-peak amplitude.
Quadratic model в 0.5.44.0 является только диагностикой и не применяется автоматически к
фотометрии.

В pipeline режим включается параметром
\texttt{--enable-psf-field-diagnostics}. По умолчанию используются только достаточно яркие,
ненасыщенные, удалённые от края и изолированные detections. Source-level результаты записываются
в \texttt{psf\_sources.csv}, а компактные frame-level coefficients --- в
\texttt{psf\_frames.csv}. Дополнительно создаются \texttt{psf\_summary.json}, карта residual PSF
по полю и временной график spatial coefficients.

Отдельная команда \texttt{psf-field-diagnostics} позволяет выполнить тот же анализ для уже
проиндексированной серии без повторного запуска aperture/ensemble photometry. Это основной
режим для ретроспективной диагностики существующих архивных данных.

Важно отличать качество изображения от пригодности фотометрии. Большой FWHM или высокая
эллиптичность сами по себе не означают, что кадр следует удалить: последующие версии
\program{} должны уметь использовать adaptive elliptical/elongated aperture и PSF/model fitting,
если поток остаётся измеримым. Версия 0.5.44.0 только измеряет морфологию и spatial structure,
чтобы дальнейшие correction policies опирались на наблюдаемую инструментальную причину.

\section{Adaptive elliptical aperture diagnostics}
\label{sec:adaptive-aperture-diagnostics}

Начиная с версии 0.5.44.1 Grandgousier умеет параллельно с базовой circular aperture
фотометрией вычислять экспериментальную adaptive elliptical aperture. В этой версии результат
остаётся диагностическим artifact и не записывается в таблицу \texttt{measurements}; следовательно,
он не может незаметно заменить science calibration. Цель этапа --- проверить переносимую
геометрию апертуры на реальных широкопольных данных прежде, чем вводить постоянное хранение
нескольких flux estimators.

Локальная форма PSF интерполируется по ярким изолированным звёздам. Для каждой reference star
используется матрица вторых моментов из раздела~\ref{sec:psf-field-diagnostics}. Для произвольной
позиции $(x,y)$ берутся ближайшие PSF references и выполняется inverse-distance weighted
интерполяция компонент covariance matrix. Из интерполированной матрицы снова находятся
${\rm FWHM}_{\rm major}$, ${\rm FWHM}_{\rm minor}$ и угол $\theta$.

Эквивалентный радиус adaptive aperture определяется как
\begin{equation}
 r_{\rm eq}=\max\left(r_0,\;k_{\rm FWHM}{\rm FWHM}_{\rm geom}\right),
\end{equation}
после чего ограничивается safety-порогом
$r_{\rm eq}\le r_0 s_{\rm max}$. Здесь $r_0$ --- базовый circular radius, а по умолчанию
$k_{\rm FWHM}=1$. Пусть
\begin{equation}
 q=\frac{{\rm FWHM}_{\rm major}}{{\rm FWHM}_{\rm minor}}.
\end{equation}
После ограничения $q$ параметром \texttt{--max-axis-ratio} полуоси равны
\begin{equation}
 a=r_{\rm eq}\sqrt{q},\qquad b=\frac{r_{\rm eq}}{\sqrt{q}}.
\end{equation}
Таким образом, апертура одновременно увеличивается при локальном уширении PSF и вытягивается
вдоль наблюдаемой главной оси. Площадь эллипса равна площади круговой апертуры радиуса
$r_{\rm eq}$.

Для вытянутых изображений background annulus автоматически сдвигается наружу, если большая
полуось приближается к исходной внутренней границе annulus. Это позволяет измерять умеренно
trailed/elongated stellar images без немедленного исключения кадра. При этом существуют safety
clamps по axis ratio и размеру; их срабатывание сохраняется в диагностических флагах.

Команда \texttt{adaptive-aperture-diagnostics} вычисляет на одних и тех же catalog-frame links
две оценки потока: существующую circular aperture и adaptive ellipse. Экспортируются
\texttt{adaptive\_aperture\_measurements.csv}, frame-level summary,
\texttt{adaptive\_aperture\_objects.csv}, JSON summary и диагностические PNG. Для object-level
сравнения из обеих instrumental magnitude series удаляются constant object median и robust
frame common mode; после этого сравниваются RMS, MAD и robust peak-to-peak. Это не заменяет
полную ensemble calibration, но позволяет проверить, уменьшает ли shape-aware aperture
пространственно-зависимую ложную переменность.

Ключевой safety-принцип версии 0.5.44.1: circular estimator остаётся эталонным science path,
а adaptive estimator применяется только как A/B диагностика. Переход к постоянному хранению
AP/adaptive-AP/PSF estimators и выбору \texttt{primary\_method} запланирован на последующие
версии после проверки на MASTER и других инструментах.

\section{Adaptive aperture как альтернативный science estimator}
\label{sec:adaptive-aperture-stage2}

Начиная с версии 0.5.44.2 adaptive elliptical aperture может использоваться не только для
диагностического CSV-сравнения, но и как отдельный источник базовой фотометрии полного Stage~2.
Режим выбирается параметром
\texttt{--photometry-estimator adaptive\_ellipse}. По умолчанию по-прежнему используется
\texttt{circular}; поэтому существующий workflow и имена circular calibration generations не
изменяются.

При adaptive-режиме source detection является обязательным внутренним prerequisite: из ярких
изолированных detections строится local PSF covariance field. Для каждого forced catalog source
интерполируется геометрия из раздела~\ref{sec:adaptive-aperture-diagnostics}, после чего поток
измеряется в ориентированном эллипсе. Если надёжный local PSF model построить нельзя, по умолчанию
используется circular fallback. Такой fallback не является скрытым: он отмечается в measurement
flags. Для строгого A/B fallback можно отключить параметром
\texttt{--no-adaptive-circular-fallback}.

Adaptive measurement использует ту же модель случайной ошибки, что и circular aperture:
source Poisson noise, background variance, неопределённость оценки background и опциональные
fractional/scintillation terms. Для эллипса число пикселей $N_{\rm ap}$ и размер background annulus
вычисляются из фактической геометрии. Поэтому
\begin{equation}
 \sigma_F^2 = \sigma_{\rm source}^2 + N_{\rm ap}\sigma_{\rm sky}^2
 + \sigma_{\rm sky\,estimate}^2 + \sigma_{\rm frac}^2 + \sigma_{\rm scint}^2
\end{equation}
применяется к реальной площади adaptive aperture, а не к площади исходного круга.

Существующая схема БД в 0.5.44.2 не расширяется. Поле
\texttt{measurements.aperture\_radius\_pix} хранит радиус круга той же площади
$r_{\rm eq}=\sqrt{ab}$, а реальные $a$, $b$, $\theta$, local FWHM, ellipticity, PSF-neighbor
coverage и границы dynamic annulus сохраняются в JSON flags. Это обеспечивает совместимость
существующих readers при полном provenance формы апертуры.

Adaptive BASE photometry получает отдельную immutable calibration generation. Далее над ней без
изменения алгоритмов выполняются обычные этапы:
\begin{equation}
\begin{aligned}
 \text{adaptive BASE}&\rightarrow\text{aperture correction}\rightarrow\text{aperture ZP}
 \rightarrow\text{ensemble}\rightarrow\text{object QC}\rightarrow{}\\
 &\qquad\text{candidate selection}\rightarrow\text{science materialization}.
\end{aligned}
\end{equation}
Таким образом, circular и adaptive варианты можно сравнивать после одной и той же downstream
calibration chain, а не только на уровне instrumental magnitudes.

Для контролируемого A/B следует задавать отдельный calibration prefix и output directory. Например:
\begin{lstlisting}
python -m grandgousier.cli process-archive-path `
  --path "C:\master\gpx\stage2\20180907_1704" `
  --execute `
  --instrument-profile MASTER `
  --photometry-estimator adaptive_ellipse `
  --calibration-prefix-template "ab_adaptive_{instrument_lower}_{date_yyyymmdd}_{field_name}_{tube}_{camera}_{filter_name}_s{series_id}_v1" `
  --out-root "exports_ab_adaptive"
\end{lstlisting}
Остальные science/QC параметры в A/B должны совпадать с circular baseline. Начиная с
версии 0.5.44.3 итоговая таблица \texttt{process-archive-path} явно печатает строку
\texttt{photometry\_estimator} для каждой pipeline-серии. Перед интерпретацией A/B результата
следует убедиться, что для adaptive-прогона там указано \texttt{adaptive\_ellipse}; это
быстрая проверка фактически применённого estimator и его provenance. Само использование
adaptive estimator не является основанием автоматически считать кадр хорошим или плохим:
морфология изображения и пригодность фотометрии остаются раздельными понятиями.

\section{Пространственно меняющаяся PSF-фотометрия}
\label{sec:spatial-psf-photometry}

Начиная с версии 0.5.44.4 полный Stage~2 поддерживает третий экспериментальный базовый
estimator: \texttt{--photometry-estimator spatial\_psf}. Значение \texttt{circular} остаётся
стандартным, а \texttt{adaptive\_ellipse} --- проверенным альтернативным aperture estimator.
Режим \texttt{spatial\_psf} предназначен прежде всего для широких полей, сложной формы PSF,
умеренного elongation и областей, где фиксированная или adaptive aperture всё ещё оставляет
пространственно-временную систематику.

Как и adaptive aperture, PSF estimator использует яркие изолированные detections для построения
локального covariance field. Для forced source в точке $(x,y)$ интерполируются
$\mathrm{FWHM}_{\rm major}$, $\mathrm{FWHM}_{\rm minor}$ и угол $\theta$. В первой реализации
используется эллиптический профиль Moffat
\begin{equation}
 P(u,v)=\frac{\beta-1}{\pi\alpha_a\alpha_b}
 \left[1+\left(\frac{u}{\alpha_a}\right)^2+
 \left(\frac{v}{\alpha_b}\right)^2\right]^{-\beta},
\end{equation}
где $u,v$ --- координаты вдоль главных осей локальной PSF. Параметры $\alpha_a,\alpha_b$
получаются из локальных FWHM при фиксированном $\beta$; стандартное значение
\texttt{--psf-moffat-beta} равно 2.5. Для диагностических сравнений доступен также
\texttt{--psf-profile gaussian}.

Положение звезды не освобождается полностью: вокруг forced catalog position выполняется небольшой
двумерный grid search centroid. Для каждой trial position поток $F$ оценивается matched-profile
least squares
\begin{equation}
 F=\frac{\sum_k P_k(D_k-B)}{\sum_k P_k^2},
\end{equation}
где $B$ --- robust median background из annulus. Центр выбирается до residual clipping по полной
PSF core; это не позволяет плохому centroid искусственно выиграть за счёт исключения центральных
пикселей. После выбора центра допускается robust clipping отдельных загрязнённых wing pixels,
причём core PSF защищён от clipping.

Для background noise используется noise-equivalent area matched filter
\begin{equation}
 N_{\rm eff}=\frac{1}{\sum_k P_k^2}.
\end{equation}
В существующую physical random-error decomposition вместо геометрического $N_{\rm ap}$
подставляется $N_{\rm eff}$. Source Poisson, uncertainty background estimate, fractional floor и
scintillation сохраняют прежнюю семантику. В поле
\texttt{measurements.aperture\_radius\_pix} для совместимости записывается
$r_{\rm noise}=\sqrt{N_{\rm eff}/\pi}$; фактическая PSF model, fit radius, centroid offsets,
residual MAD, capture fraction и $N_{\rm eff}$ сохраняются в flags. Centroid offsets также
заполняют нормализованные поля \texttt{centroid\_dx\_pix}/\texttt{centroid\_dy\_pix}.

Если профиль конкретного источника не удаётся измерить, стандартный fallback ---
adaptive ellipse с той же локальной геометрией. Его можно отключить параметром
\texttt{--no-psf-adaptive-fallback}. Если невозможно построить само PSF field, последним
совместимым fallback остаётся circular aperture, если не указан
\texttt{--no-adaptive-circular-fallback}. Каждый fallback явно сохраняется в provenance.

PSF BASE является отдельной immutable calibration generation. Обычный aperture-correction этап
при этом не считается логически лишним: для PSF photometry он играет роль стандартной
profile-to-reference-aperture correction и компенсирует возможную ошибку нормировки wings или
несовпадение аналитического профиля с реальной PSF. После этого используются те же aperture-ZP,
ensemble, object-QC, candidate selection и science materialization, что и для circular/adaptive
A/B.

Первый контрольный тест выполняется на синхронных MASTER series 76 EAST/V и 77 WEST/R. Circular
baseline дал 9/143 финальных candidates, а полный adaptive Stage~2 --- 7/30 при одновременном
росте accepted science points. Цель spatial-PSF A/B --- проверить, уменьшится ли остаточная
сильно коррелированная R candidate population без ухудшения V и без потери большого числа
science measurements.

\section{Эмпирическая пространственная PSF-фотометрия}
\label{sec:empirical-psf-photometry}

Реальный A/B-тест аналитического \texttt{spatial\_psf} на MASTER GPX\_TF17 показал, что
эллиптический Moffat с фиксированным $\beta=2.5$ не является достаточно точной моделью широкого
поля WEST/R. По сравнению с проверенным \texttt{adaptive\_ellipse} число финальных R-candidates
изменилось с 30 до 144, accepted science points уменьшились с 29312 до 26637, а median
RMS/MAD ухудшились с 0.05240/0.01302 до 0.06118/0.02167 mag. Характерная
profile-to-reference aperture correction около $+0.46\ldots+0.47$ mag показывает, что аналитическая
нормировка wings не соответствует реальной PSF. Поэтому начиная с 0.5.44.5 добавлен отдельный
экспериментальный estimator
\begin{lstlisting}
--photometry-estimator empirical_psf
\end{lstlisting}
который не предполагает аналитическую форму PSF.

На каждом кадре выбираются яркие, изолированные и ненасыщенные detections. Для каждой reference
star вычисляется субпиксельный centroid, после чего изображение звезды пересэмплируется на общую
сетку offsets $(\Delta x,\Delta y)$ билинейной интерполяцией. Из stamp вычитается robust local
background. Затем stamp нормируется не к бесконечному total flux, а к потоку внутри большой
reference aperture радиуса $r_{\rm ref}$:
\begin{equation}
 P_j(\Delta x,\Delta y)=\frac{I_j(\Delta x,\Delta y)-B_j}
 {\sum_{r\le r_{\rm ref}}[I_j-B_j]}.
\end{equation}
По умолчанию $r_{\rm ref}=5$ pix, то есть совпадает со стандартным reference radius downstream
aperture-correction слоя.

В поле detector создаётся регулярная сетка PSF nodes, по умолчанию $3\times3$. Каждый node строится
из ближайших reference stars, а не только из источников внутри жёсткой прямоугольной ячейки. Это
уменьшает риск пустых corner nodes. Для каждого pixel узловой template определяется robust
median/clip комбинацией зарегистрированных нормированных stamps. Между четырьмя соседними nodes
интерполируется уже полный двумерный template
\begin{equation}
 P(x,y;\Delta x,\Delta y),
\end{equation}
а не только FWHM, ellipticity или orientation.

Для forced source сохраняется небольшой centroid grid search. При заданном локальном template
поток оценивается matched-profile выражением
\begin{equation}
 F=\frac{\sum_k P_k(D_k-B)}{\sum_k P_k^2}.
\end{equation}
Так как template нормирован к reference aperture, $F$ непосредственно оценивает поток на той же
фотометрической шкале. Поэтому downstream aperture correction должна быть мала и в основном
диагностировать residual mismatch, а не компенсировать систематическую ошибку порядка 0.5 mag,
как в неудачном аналитическом Moffat A/B.

Для random-error model используется noise-equivalent area
\begin{equation}
 N_{\rm eff}=\frac{1}{\sum_k P_k^2},
\end{equation}
после чего применяются те же source Poisson, measured/detector background, uncertainty background
estimate, fractional floor и scintillation terms, что и для других BASE estimators.

По умолчанию fallback hierarchy имеет вид
\begin{equation}
 \text{empirical PSF}\rightarrow\text{adaptive ellipse}\rightarrow\text{circular aperture}.
\end{equation}
Первый переход можно запретить параметром
\texttt{--no-empirical-psf-adaptive-fallback}, последний --- существующим
\texttt{--no-adaptive-circular-fallback}. Каждый fallback сохраняется в measurement flags.

Основные empirical параметры:
\begin{lstlisting}
--empirical-psf-grid-size 3
--empirical-psf-stamp-radius-pix 8
--empirical-psf-reference-radius-pix 5.0
--empirical-psf-references-per-node 48
--empirical-psf-min-references-per-node 12
\end{lstlisting}
Centroid и residual clipping используют общие PSF controls
\texttt{--psf-centroid-search-radius-pix}, \texttt{--psf-centroid-grid-steps},
\texttt{--psf-residual-clip-sigma} и \texttt{--psf-residual-clip-iterations}.

Для MASTER acceptance gate empirical PSF сравнивается прежде всего с adaptive baseline:
7/30 финальных V/R candidates, 30995/29312 accepted points и R median RMS/MAD
0.05240/0.01302 mag. Простое превосходство circular estimator недостаточно: новый estimator должен
как минимум не ухудшать adaptive result и желательно уменьшать остаточную detector-Y/common-mode
структуру 30 R-candidates.


\section{Пространственная ensemble-калибровка}
\label{sec:spatial-ensemble}

Начиная с версии 0.5.44.8 внутренний conventional-ensemble base перед пространственной коррекцией обязательно материализуется и публикуется как отдельная immutable generation с полным набором измерений. Это устраняет неоднозначность при повторном использовании calibration prefix: spatial stage всегда читает только что построенную base-generation, а не старый normalized-only head. Финальный spatial product имеет provenance suffix \texttt{spatial\_v2}; пользовательские параметры и математическая модель при этом не изменились.


После реальных A/B-тестов MASTER предпочтительным BASE estimator остаётся
\texttt{adaptive\_ellipse}. Аналитический и empirical PSF fitting не превзошли его, поэтому
начиная с 0.5.44.6 остаточная широкопольная систематика корректируется на уровне ensemble, а не
путём дальнейшего усложнения модели отдельной звезды.

Обычная ensemble-калибровка сначала определяет reference stars и frame-common offset $Z_j$.
Затем для reference residuals оценивается
\begin{equation}
 m_{ij}=M_i+Z_j+S_j(x_i,y_i)+\epsilon_{ij}.
\end{equation}
Здесь $S_j$ --- только дифференциальная пространственная часть. В fit присутствует nuisance
intercept, однако он не применяется к science measurements. Пространственные basis functions
центрируются по reference stars каждого кадра, поэтому weighted mean применяемой поправки равен
нулю и новый слой не дублирует обычный frame zero point.

Рассматриваются identity, plane и quadratic модели. Выбор выполняется по BIC, минимальной
spatial peak-to-peak amplitude и safety limit на максимальную поправку по полю. Используется
iterative robust clipping. Начиная с 0.5.44.9 добавлен обязательный по умолчанию robust
non-degradation gate: предложенная plane/quadratic модель применяется только если её residual MAD
не превышает residual MAD identity-модели более чем на
\texttt{--spatial-ensemble-max-residual-mad-degradation}. Default равен 0 mag, то есть увеличение
MAD запрещено. Отрицательное значение CLI отключает этот gate для диагностического воспроизведения
старого поведения. Провал этого локального gate не делает кадр плохим и не уменьшает
\texttt{good\_frame\_fraction}: кадр остаётся валидным, но получает identity correction. При
недостатке reference stars или провале frame quality также применяется identity. Если доля пригодных
кадров ниже global threshold, вся final spatial calibration становится identity-copy исходного
ensemble product.

В 0.5.44.6 random-error columns намеренно не увеличиваются uncertainty spatial model: uncertainty
сохраняется только в flags. Это делает A/B проверку консервативной --- уменьшение candidate
population должно происходить за счёт исправления light curves, а не за счёт увеличения
$\sigma_m$.

Для включения:
\begin{lstlisting}
--enable-spatial-ensemble
--spatial-ensemble-max-order 2
--spatial-ensemble-min-ref-stars-per-frame 30
--spatial-ensemble-min-bic-improvement 6
--spatial-ensemble-min-spatial-peak-to-peak 0.003
--spatial-ensemble-max-residual-mad-degradation 0.0
--spatial-ensemble-max-abs-spatial-deviation 0.15
\end{lstlisting}
При candidate-excluded ensemble iterations spatial surface строится заново после каждого нового
набора reference stars. Диагностика записывается в каталог \texttt{spatial\_ensemble/}.

Диагностика сохраняется раздельно по ensemble-итерациям в каталогах
\texttt{spatial\_ensemble/iter1}, \texttt{iter2}, ... . Это важно при
candidate-excluded convergence: результаты предыдущей spatial-модели не
перезаписываются последующими итерациями. В 0.5.44.9 frame CSV/JSON дополнительно сохраняют
\texttt{proposed\_model}, proposed/selected $\Delta\mathrm{MAD}$, fractional MAD improvement,
$\Delta\mathrm{RMS}$ и \texttt{fallback\_reason}. Это позволяет отличить обычный identity выбор
от spatial-модели, отвергнутой именно robust non-degradation gate. Новая spatial generation имеет
suffix \texttt{spatial\_v3} и algorithm provenance
\texttt{robust\_zero\_mean\_spatial\_residual\_v3}.


Начиная с версии 0.5.44.10 spatial-ensemble продукты, создаваемые основным pipeline,
подчиняются общему механизму неизменяемых поколений калибровок. Поэтому повторный запуск
с тем же логическим именем безопасен: завершённый продукт получает новое физическое поколение,
а незавершённый продукт от прерванного запуска может быть повторно материализован до публикации.
Это изменение не затрагивает математическую модель spatial correction.Начиная с версии 0.5.44.11 spatial-ensemble слой также записывает компактное нормализованное
представление своей модели. Для каждого кадра сохраняется \texttt{CalibrationFrameTerm} с
типом \texttt{SPATIAL\_ENSEMBLE}: константа и линейные коэффициенты используют штатные поля
frame term, quadratic коэффициенты и границы детектора хранятся в существующем JSON
\texttt{details}/parameters. Это не изменяет science correction и не требует миграции схемы.
Dual-write нужен для immutable calibration reconstruction: downstream-код может воспроизвести
пространственную поправку как
\begin{equation}
 m_{\mathrm{out}} = m_{\mathrm{source}} - S_j(x_n,y_n),
\end{equation}
при этом \texttt{mag\_err}, \texttt{mag\_err\_random} и \texttt{mag\_err\_systematic} остаются
неизменными. Полные \texttt{Measurement} rows пока сохраняются для контролируемого A/B, поэтому
0.5.44.11 является исправлением provenance/reconstruction contract, а не новой научной моделью.

\section{Консервативный post-ensemble photometric frame QC}
\label{sec:photometric-frame-qc}

Начиная с версии 0.5.44.12 доступен режим
\begin{lstlisting}
--enable-photometric-frame-qc
\end{lstlisting}
Его задача --- не исправлять неизвестную инструментальную систематику сложной моделью, а автоматически
оценивать надёжность фотометрии каждого кадра после обычной ensemble-калибровки. Режим взаимоисключающий
с \texttt{--enable-spatial-ensemble}: spatial ensemble изменяет величины, тогда как photometric frame QC
использует пространственную модель только как диагностический признак.

Для reference stars обычного ensemble определяется residual
\begin{equation}
 r_{ij}=m_{ij}-\widetilde m_i,
\end{equation}
где $\widetilde m_i$ --- robust центральная величина reference star по серии. Для каждого кадра
вычисляется
\begin{equation}
 R_j=1.4826\,\operatorname{MAD}(r_{ij}),
\end{equation}
а характерный уровень серии
\begin{equation}
 R_0=\operatorname{median}_j R_j.
\end{equation}

\subsection{Политика 0.5.44.14: QC не подменяет модель ошибки}

Реальный A/B 0.5.44.13 подтвердил, что hard-reject-only классификация сохраняет значительно больше
эпох: EAST/V получил 17 GOOD, 6 SUSPECT и 2 BAD кадра, WEST/R --- 13 GOOD, 10 SUSPECT и 2 BAD.
Accepted science points восстановились до 29964 и 29589. Однако добавление одной и той же
frame-specific величины ошибки всем звёздам SUSPECT-кадра увеличило final candidates до 39 и 49
при adaptive-only baseline 7 и 30. Поэтому начиная с 0.5.44.14 QC не изменяет ошибки GOOD/SUSPECT
измерений.

Каждый кадр сохраняет один из статусов:
\begin{description}
  \item[GOOD] действие \texttt{KEEP}; величина и все компоненты ошибки сохраняются;
  \item[SUSPECT] действие \texttt{KEEP\_FLAGGED}; magnitude и
  \texttt{mag\_err}, \texttt{mag\_err\_random}, \texttt{mag\_err\_systematic} не изменяются;
  \item[BAD] действие \texttt{REJECT}; исходная фотометрия остаётся в базе, но производная QC-калибровка
  выставляет \texttt{mag\_cal=NULL}. BAD разрешён только для hard failures.
\end{description}

К hard failures относятся невозможность получить residual scatter, катастрофически большой абсолютный
robust scatter (в текущей policy $R_j\ge0.10$ mag) и критически малое число reference stars
(менее пяти). Relative outlier внутри серии, недостаточное, но не критическое число references и
пространственная residual structure сами по себе дают максимум SUSPECT.

\subsection{Frame-specific scatter как диагностическая величина}

По-прежнему вычисляется
\begin{equation}
 \sigma_{\mathrm{frame},j}
 =\sqrt{\max\left(R_j^2-R_0^2,0\right)}.
 \label{eq:frame-qc-uncertainty}
\end{equation}
Но эта величина является характеристикой residual population кадра, а не доказанной одинаковой
неопределённостью каждой отдельной звезды. Поэтому в 0.5.44.14 она сохраняется только как
\texttt{photometric\_frame\_qc\_sigma\_frame\_diagnostic\_mag} и не складывается ни с одной
measurement-error component. Такая политика избегает статистической подгонки коэффициента error inflation.

\subsection{Подтверждение кандидатов по GOOD-эпохам}

Для автоматического candidate search используется уже существующий field-QC/candidate evaluator, но
после обычного анализа всех пригодных эпох он повторно запускается на подмножестве кадров со статусом
GOOD. Новая статистическая модель не вводится: используются те же RMS/MAD/intrinsic criteria и те же
пороги.

Автоматический кандидат остаётся \texttt{ACCEPTED}, только если он принят в обоих анализах:
\begin{description}
  \item[GOOD\_FRAME\_CONFIRMED] accepted на всех usable epochs и на GOOD epochs;
  \item[SYSTEMATICS\_SENSITIVE] accepted только при участии SUSPECT epochs; строка переводится в
  \texttt{BORDERLINE}, но сохраняется как диагностический объект;
  \item[INSUFFICIENT\_GOOD\_FRAME\_SUPPORT] GOOD-кадров меньше минимального числа, необходимого обычному
  object/candidate evaluator; автоматический кандидат переводится в \texttt{BORDERLINE};
  \item[TARGET\_PRESERVED] явно заданный target не удаляется автоматическим confirmation gate.
\end{description}

Confirmation выполняется до candidate-excluded ensemble. Поэтому объект, который выглядит переменным
только из-за SUSPECT-эпох, не исключается из reference ensemble и не создаёт положительную обратную связь
между ложным кандидатом и следующей итерацией калибровки. Тот же подтверждённый набор используется при
финальном экспорте candidate light curves. Raw all-usable triage сохраняется в JSON provenance для A/B.

Начиная с 0.5.44.15 для уже \texttt{GOOD\_FRAME\_CONFIRMED} кандидата флаг Object QC
\texttt{HIGH\_RMS} сохраняется как диагностический признак, но больше не является причиной отклонить
все точки его candidate light curve. Это узкое исключение устраняет логический цикл ``объект подтверждён
как переменный по повышенному scatter, затем тот же повышенный scatter удаляет его кривую''. Все остальные
science reject flags сохраняются без изменений; для обычных объектов и неподтверждённых candidates
\texttt{HIGH\_RMS} по-прежнему действует как штатный reject flag. В candidate export явно записывается
\texttt{HIGH\_RMS\_DIAGNOSTIC\_ONLY\_FOR\_GOOD\_FRAME\_CONFIRMED}.

CSV \texttt{field\_qc\_candidate\_triage.csv} также сохраняет статус GOOD-frame confirmation и
GOOD-only RMS/MAD/robust diagnostics, поэтому distinction между подтверждённой и
systematics-sensitive переменностью доступно без разбора JSON provenance.


\subsection{Корреляция кандидатов внутри полос детектора (0.5.44.16--0.5.44.17)}

После \texttt{GOOD\_FRAME\_CONFIRMED} Grandgousier дополнительно выполняет диагностическую проверку,
не меняющую статус кандидата. Цель --- заметить случай, когда несколько кандидатов в одной полосе CCD
синхронно меняют блеск даже на кадрах, которые глобальный frame QC считает GOOD.

Для каждого accepted confirmed candidate по GOOD-эпохам определяется медианное положение
$(\widetilde x_i,\widetilde y_i)$ и центрированная величина
\begin{equation}
 \Delta m_i(t)=m_i(t)-\operatorname{median}_{t\in GOOD}m_i(t).
\end{equation}
Каждая ось детектора независимо делится на 16 bins. Чтобы близкая пара не перестала сравниваться только
из-за случайного попадания по разные стороны границы bin, параллельно используется вторая разметка тех же
16 bins, сдвинутая на половину ширины. Пара кандидатов сравнивается по X или Y, если она находится в
одном bin хотя бы одной из двух разметок соответствующей оси.

Для минимум восьми общих GOOD-эпох начиная с 0.5.44.17 параллельно вычисляются два коэффициента:
\begin{equation}
 \rho_{S,ij}=\rho_S\left(\Delta m_i(t),\Delta m_j(t)\right),
\end{equation}
\begin{equation}
 r_{P,ij}=r_P\left(\Delta m_i(t),\Delta m_j(t)\right).
\end{equation}
Spearman $\rho_S$ чувствителен к устойчивому сходству формы/рангов световых кривых, тогда как Pearson
$r_P$ дополнительно чувствителен к синхронным высокоамплитудным всплескам или провалам, даже если остальные
эпохи слабо согласованы. Пара считается диагностически коррелированной, если
\begin{equation}
 |\rho_{S,ij}|\ge0.8 \quad \text{или} \quad |r_{P,ij}|\ge0.8.
\end{equation}
Добавляется флаг соответствующей оси:
\begin{itemize}
  \item \texttt{CANDIDATE\_CORRELATED\_X};
  \item \texttt{CANDIDATE\_CORRELATED\_Y}.
\end{itemize}
Знаки обоих коэффициентов сохраняются. Антикорреляция также информативна для
blending/PSF/flux-redistribution artefacts. Для каждой пары и кандидата сохраняется режим
\texttt{SPEARMAN}, \texttt{PEARSON}, \texttt{BOTH} или \texttt{NONE}. Флаг не переводит кандидат в
BORDERLINE/REJECTED и не влияет на candidate-excluded ensemble.

В candidate triage и final candidate summary сохраняются detector X/Y, primary и half-shifted bins,
наиболее сильные signed Spearman/Pearson coefficients, object/source ID соответствующих партнёров,
число общих GOOD-эпох, число коррелирующих соседей по каждой метрике и итоговый correlation mode.
Полный список проверенных пар экспортируется в
\texttt{candidate\_axis\_correlation/candidate\_axis\_pairs.csv}, summary --- в соседний
\texttt{summary.json}.

Реальный 0.5.44.16 A/B подтвердил wiring исходной Spearman-диагностики: 0/5 confirmed V и 17/25
confirmed R candidates получили axis-correlation flag при неизменных science photometry и candidate
acceptance. Разбор тех же кривых показал, что общий одиночный всплеск может иметь слабый Spearman, но
очень сильный Pearson. Поэтому 0.5.44.17 расширяет только диагностическую метрику; фотометрия, frame QC,
GOOD-frame confirmation и candidate acceptance не меняются.

\subsection{Spatial diagnostics без изменения magnitude}

Существующий robust fitter residual surface plane/quadratic переиспользуется только для обнаружения
\texttt{SPATIAL\_SYSTEMATIC}. Найденная поверхность $S_j(x,y)$ не вычитается из magnitude. Таким
образом изменение механики, юстировки, камеры или оптики не требует заранее известной epoch-dependent
instrument map и не создаёт произвольную frame-to-frame spatial correction.

\subsection{PSF-параметры в QC}

Повторное чтение FITS не выполняется. QC объединяет два уже существующих источника:
aperture/PSF context и покадровый экспорт \texttt{psf\_frames.csv} команды PSF field diagnostics.
В \texttt{frames.quality\_flags} сохраняется компактный \texttt{frame\_psf\_summary\_v1}, содержащий
число PSF sources, median geometric FWHM (pix и, когда доступно, arcsec), major/minor FWHM,
ellipticity, axial position angle, а также spatial peak-to-peak для FWHM и ellipticity. Начиная с
0.5.44.15 dedicated \texttt{psf\_frames.csv} читается по его реальным полям
\texttt{fwhm\_geom\_*} и \texttt{ellipticity\_*}; отдельно для FWHM и ellipticity выбирается успешная
constant/plane/quadratic модель с минимальным BIC. В summary сохраняются и имя выбранной модели, и её
spatial peak-to-peak. Эти величины остаются исключительно диагностическими.

PSF сам по себе не является reject criterion: необычная форма PSF может быть корректно обработана
\texttt{adaptive\_ellipse}. PSF-параметры нужны для физической интерпретации причин деградации и
последующего анализа архива.

\subsection{Хранение и provenance}

Photometric QC записывается как normalized immutable stage \texttt{CalibrationIteration};
тип итерации равен \texttt{PHOTOMETRIC\_FRAME\_QC}. В 0.5.44.14 новая logical generation имеет
suffix \texttt{photqc\_v3}; алгоритм называется
\texttt{post\_ensemble\_photometric\_frame\_qc\_v3}. Совместимые блоки
\texttt{frame\_photometric\_qc\_v1} и \texttt{frame\_psf\_summary\_v1} остаются в
\texttt{frames.quality\_flags}. Миграция БД не требуется.

Режим остаётся opt-in до повторного MASTER A/B. Production acceptance требует одновременно сохранить
почти весь набор не-BAD эпох и существенно уменьшить число автоматических candidates, зависящих от
SUSPECT-кадров, без изменения photometric error model.
